\section{Powersets}

\begin{defn}
    If $A$ is a set, the \emph{power set} of $A$ is another set, denoted $\powerset(A)$ and defined to be the set of \emph{all} subsets of $A$. Formally, 
    $$
    \powerset(A)=\{ X:X\subseteq A \}.
    $$
\end{defn}

Suppose that we have some set $A=\{ 1,2,3 \}$, then 
$$\powerset(A)=\{ \varnothing,\{ 1 \},\{ 2 \},\{ 3 \}, \{ 1,2 \},\{ 1,3 \},\{ 2,3 \},\{ 1,2,3 \} \}.$$ 
A \emph{finite} set of $n$-elements has $2^n$ subsets and the powerset of that set has $2^n$-elements. 

\begin{fact}
    If $A$ is a finite set, then $\lvert \powerset(A) \rvert=2^{\lvert A \rvert}$.
\end{fact}

\begin{example}
    $$
    \powerset(\{ a \})\times\powerset(\{ \varnothing \})=\{ (\{ a \},\{ \varnothing \}), (\{ a \},\varnothing), (\varnothing,\varnothing),(\varnothing, \{ \varnothing \}) \}.
    $$
    This example might appear a little bit intimidating, however consider the following: 
    $$
    \powerset(\{ a \})=\{ \varnothing,\{ a \} \} \quad \powerset(\{ \varnothing \})=\{ \varnothing,\{ \varnothing \} \}.
    $$
    Then we simply pair the elements together. 
\end{example}

Next, some \emph{wrong} facts about powersets. 
\begin{enumerate}
    \item $\powerset(1)=\{ \varnothing,\{ 1 \} \}$ is meaningless, since $1$ is not a set.\footnote{There could be some debate here, this is what is says in the textbook, so I will follow it, however we can express $1=\{ \varnothing \}$ or $\{ \varnothing, \{ \varnothing \} \}$ when we consider the $V$-hierarchy of set construction.}
    \item $\powerset(\{ 1,\{ 1,2 \} \})=\{ \varnothing,\{ 1 \},\{ \{ 1,2 \} \},\{ 1,\{ 1,2 \} \} \}$ becuase $\{ 1,2 \}\not\subseteq \{ 1,\{ 1,2 \} \}$, however $\{ 1,2 \}\in\{\{ 1,\{ 1,2 \}\}\}$.\footnote{The correct statement would be $\{ \{ 1,2 \} \}\in\powerset(\{ 1,\{ 1,2 \} \})$.}
    \item $\powerset(\{ 1,\{ 1,2 \} \})$ but $\{ \{ 1 \} \}\not\subseteq \{ 1,\{ 1,2 \} \}$, the reasoing is similar as in the previous footnote. 
\end{enumerate}

It is not possible to list all the members of a powerset of $A$ if $A$ is \emph{infinite.}

\subsection{Exercises From Section 4}
\begin{exercise}\label{ex:4.1}
    Write the following sets by listing their elements between braces. 
    \begin{enumerate}
        \item $\powerset(\{ \{ a,b \},\{ c \} \})$
        \item $\powerset(\{ 1,2,3,4 \})$
        \item $\powerset(\{ \{ \varnothing \},5 \})$
        \item $\powerset(\{ \R,\Q \})$
        \item $\powerset(\powerset(\{ 2 \}))$
        \item $\powerset(\{ 1,2 \})\times\powerset(\{ 3 \})$
        \item $\powerset(\{ a,b \})\times\powerset(\{ 0,1 \})$
        \item $\powerset(\{ 1,2 \}\times \{ 3 \})$
        \item $\powerset(\{ a,b \}\times \{ 0 \})$
        \item $\{ X\in\powerset(\{ 1,2,3 \}):\lvert X \rvert\le 1 \}$
        \item $\{ X\subseteq\powerset(\{ 1,2,3 \}):\lvert X \rvert\le 1 \}$
        \item $\{ X\in\powerset(\{ 1,2,3 \}):2\in X \}$.
    \end{enumerate}
\end{exercise}\hyperref[sol:4.1]{[Go to Solution]}

\begin{exercise}\label{ex:4.2}
    Suppose that $\lvert A \rvert=m,\lvert B \rvert=n$. Find the following cardinalities. 
    \begin{enumerate}
        \item $\lvert \powerset\powerset\powerset(A) \rvert$
        \item $\lvert \powerset\powerset(A) \rvert$
        \item $\lvert \powerset(A\times B) \rvert$
        \item $\lvert \powerset(A)\times\powerset(B) \rvert$
        \item $\lvert \{ X\in\powerset(A):\lvert X \rvert\le 1 \} \rvert$
        \item $\lvert \powerset(A\times\powerset(B)) \rvert$
        \item $\lvert \powerset\powerset\powerset(A\times\varnothing) \rvert$
        \item $\lvert \{ X\subseteq\powerset(A):\lvert X \rvert\le 1 \} \rvert$.
    \end{enumerate}
\end{exercise}\hyperref[sol:4.2]{[Go to Solution]}
